\section{Introduction}

Spatial transcriptomics is an emerging class of technologies that couples gene expression measurement with the spatial organisation of tissue, allowing researchers to relate transcriptional state to histological context. These technologies are commonly grouped into two families. Sequencing-based methods, such as the original Spatial Transcriptomics platform~\cite{stahl2016spatialtranscriptomics} and its commercial successors including Visium, deposit barcoded spots on a slide and read transcripts with short-read sequencing. Imaging-based methods, such as MERFISH~\cite{chen2015merfish,moffitt2018merfish}, seqFISH+~\cite{eng2019transcriptome} and the commercial Xenium platform~\cite{janesick2023xenium,rao2021xenium}, localise a targeted panel of transcripts through successive rounds of fluorescent hybridisation. Sequencing-based methods are attractive because they provide transcriptome-wide measurements, recover splice variants, and permit downstream analyses such as RNA velocity~\cite{lamanno2018rnavelocity}, but until recently they have lacked the spatial resolution required to resolve individual cells.

Visium HD~\cite{10xvisiumhd2024} was introduced to close this resolution gap. The platform collects transcript counts at $2\,\mu\mathrm{m}\times 2\,\mu\mathrm{m}$ bins, well below the size of a mammalian cell, and exposes an aggregated $8\,\mu\mathrm{m}\times 8\,\mu\mathrm{m}$ view that approximates cellular dimensions. In practice, however, neither representation aligns cleanly to cells. Smaller cells are contaminated by transcripts from their neighbours when their profile is derived from an $8\,\mu\mathrm{m}$ spot, while cells that are either small or tightly packed commonly straddle multiple aggregated spots. Recovering accurate single-cell transcript estimates from Visium HD therefore requires a mapping step between bin-level measurements and segmented cell boundaries.

Prior work tackles this mapping with imaging. Bin2cell~\cite{polacco2024bin2cell} pairs the Stardist~\cite{schmidt2018stardist} segmentation model---a star-convex extension of U-Net~\cite{ronneberger2015unet}---with Visium HD to produce cell outlines, dilates those outlines, and then assigns $2\,\mu\mathrm{m}$ bins to a cell when the bin centroid falls within the expanded outline. This naive, unique-assignment strategy works well for sparse tissues but loses information whenever bins overlap several cells simultaneously. In dense epithelial or stromal regions, such overlaps dominate: our evaluations show that up to one in four bins overlap more than one cell. Methods that distribute transcripts among all overlapping cells, rather than discarding the ambiguous bins, are therefore necessary to obtain accurate per-cell expression in the dense regime.

To address this gap we developed ENACT, an end-to-end, tissue-agnostic pipeline for Visium HD. ENACT couples Stardist segmentation with four alternative bin-to-cell assignment strategies---one naive and three weight-based schemes---and with three established cell-type annotation engines: the probabilistic CellAssign model~\cite{zhang2019probabilistic}, the logistic-regression-based CellTypist classifier~\cite{conde2022celltypist}, and the score-based Sargent annotator~\cite{nouri2024sargent} that is particularly well suited to the low transcript counts typical of spatial assays. The pipeline ingests a 10x Genomics Visium HD bin-by-gene matrix and a paired full-resolution histology image, and emits AnnData~\cite{virshup2024anndata} objects compatible with Squidpy~\cite{palla2022squidpy} for neighbourhood enrichment, Moran's I~\cite{moran1950notes} and co-occurrence analyses. Cell locations can be visualised directly in TissUUmaps~\cite{pielawski2023tissuumaps}.

Our contributions are fourfold. First, ENACT is, to our knowledge, the first tissue-agnostic end-to-end pipeline that turns raw Visium HD bin-level data into annotated, spatially resolved cell objects. Second, we introduce three weight-based bin-to-cell assignment strategies that generalise the naive scheme used by prior work and recover information otherwise lost in densely packed tissue. Third, we evaluate all four assignment strategies and all three annotation engines on synthetic Visium HD datasets derived from Xenium and seqFISH+ ground truth, and on 20{,}991 pathologist-annotated cells from human colorectal cancer and mouse small intestine, establishing the Weighted-by-Area strategy combined with Sargent as the strongest default. Fourth, the pipeline is released as open source and produces paired image and gene-expression datasets at scale, which can be used as training data for future vision-based cell classifiers.

\section{Related Work}

\paragraph{Spatial transcriptomics platforms.}
Visium HD sits between the older Visium technology---which samples multi-cell spots~\cite{stahl2016spatialtranscriptomics}---and imaging-based technologies that achieve sub-cellular resolution on a targeted panel~\cite{chen2015merfish,moffitt2018merfish,eng2019transcriptome,janesick2023xenium,rao2021xenium}. Its strength is transcriptome coverage combined with a sub-cellular grid, but the cell-level resolution has to be recovered in silico; imaging-based methods deliver cell-level resolution directly but are limited to several hundred target genes. ENACT is designed to let sequencing-based studies reach the same granularity as imaging platforms by leveraging the paired histology image.

\paragraph{Cell segmentation.}
Cell segmentation is a prerequisite for any bin-to-cell assignment. U-Net~\cite{ronneberger2015unet} set the standard for biomedical image segmentation, and Stardist~\cite{schmidt2018stardist} extended it with star-convex polygon outputs that handle densely packed nuclei particularly well. ENACT adopts Stardist by default because packed epithelia and tumour regions are the very settings where assignment ambiguity is largest.

\paragraph{Bin-to-cell assignment for Visium HD.}
Bin2cell~\cite{polacco2024bin2cell} is, to our knowledge, the only prior published method that couples Stardist segmentation with Visium HD bin assignment. It uses a naive, unique-assignment rule: bins whose centroid falls within an expanded cell outline are summed into that cell, and overlapping bins are effectively discarded. ENACT both reproduces this strategy and extends it with three weight-based alternatives that distribute overlapping transcripts among all cells the bin intersects, and we quantify the gain empirically.

\paragraph{Cell-type annotation.}
Cell-type annotation from single-cell expression has converged on three broad families: probabilistic marker-based methods such as CellAssign~\cite{zhang2019probabilistic}, machine-learning classifiers trained on large reference atlases such as CellTypist~\cite{conde2022celltypist}, and score-based marker methods such as Sargent~\cite{nouri2024sargent}. Marker lists are sourced from curated resources including PanglaoDB~\cite{franzen2019panglaodb} and CellMarker~\cite{zhang2019cellmarker}. We include all three families in ENACT to give users flexibility and to quantify how annotation choice interacts with bin-to-cell assignment.

\paragraph{Downstream spatial analysis.}
Once single-cell expression and cell types are in hand, downstream spatial statistics---neighbourhood enrichment, Moran's I, co-occurrence---are most conveniently computed with Squidpy~\cite{palla2022squidpy}, which builds on Scanpy~\cite{wolf2018scanpy} and the AnnData data model~\cite{virshup2024anndata}. ENACT writes directly into this ecosystem, and cells can be inspected in TissUUmaps~\cite{pielawski2023tissuumaps}.
